\documentclass[11pt,a4paper]{article}
\usepackage[margin=20mm]{geometry}
\usepackage{amsmath,amssymb,booktabs}
\usepackage[T1]{fontenc}
\usepackage{lmodern}
\usepackage[hidelinks]{hyperref}
\hypersetup{pdftitle={Joint PCA-QUBO margin: European and American options, version 2}}
\setlength{\parindent}{0pt}
\setlength{\parskip}{5pt}
\setlength{\emergencystretch}{2em}
\newcommand{\pnl}{\operatorname{PnL}}
\newcommand{\diag}{\operatorname{diag}}
\begin{document}
\begin{center}
{\Large Joint PCA--QUBO margin for equities,\\European and American options}\\[5pt]
{\large Version 2: implemented extension}\\[4pt]
{\small 12 September 2026}
\end{center}

\textbf{What changed.}
The original document described equities and European equity/spot-index options.
The new research API accepts vanilla calls and puts of either exercise style,
with long or short positions. European values use Black--Scholes; American values
use the Ju--Zhong approximation specified in the supplied
\path{american_pca_qubo_model.md}. Calibration, horizon valuation and sensitivities
use the same selected pricing formula. Boundary-aware candidate generation
supplements QUBO proposals, and final ranking uses full nonlinear portfolio P\&L.

The existing PCA geometry and integer-ball QUBO algebra are retained. No binary
exercise choice is added. With three stress coordinates and eight bits per
coordinate, the QUBO still has 122 variables, regardless of option count.
Multiple smooth anchors generate separate QUBOs on the same stress domain.

\textbf{Scope.}
The supported instruments are vanilla equity and spot-index options with
continuous dividend yield and European or continuous American exercise.
American pricing requires nonnegative rates and yields, positive spot, strike
and IV, and nonnegative remaining maturity. Rates, yield and IV are held fixed
across scenarios. Futures options, discrete dividends, exotic exercise/payoffs,
FX and stochastic volatility are outside this extension.

This is a frozen-book horizon revaluation. It treats each option as outstanding
at the horizon; it does not simulate intervening exercise, assignment or their
cash/stock ledger. A short American contract is a negative holding of the
holder's value, not a different exercise policy. Ju--Zhong is an approximation
to American optimal stopping. The result is the \emph{greatest model loss found},
not an exact American value, a certified worst-case loss, or a calibrated VaR.

\textbf{Implementation map.}
\begin{center}\small
\begin{tabular}{p{.43\linewidth}p{.49\linewidth}}
\toprule
Component & Responsibility\\\midrule
\texttt{VanillaPriceContext} & Immutable pricing context, cached boundary, branch derivatives and domain checks.\\
\texttt{OptionBook} & Dated marks, contract units, bracketed IV calibration, supplied-IV diagnostics and nonlinear option P\&L.\\
\texttt{OptionFactorStressModel} & Shared stock/option scenarios, residual alignment and mixed quadratic coefficients.\\
\texttt{solveOptionFactorStress} & Boundary seeds, smooth-anchor QUBOs, integer repair and nonlinear ranking.\\
\bottomrule
\end{tabular}\end{center}
The original European-only API and PDF remain available. Application YAML,
existing option-market conventions and T0--T3 runner defaults are unchanged.

\newpage
\textbf{1. Portfolio units, dates and shared scenarios}

For underlying $i$, let $S_{i0}>0$ be the dated spot and $w_i$ its signed monetary
return exposure. A holding of $q_i^{\rm sh}$ shares has $w_i=q_i^{\rm sh}S_{i0}$.
For option $o$, let $i(o)$ identify its underlying, $K_o$ its strike, $T_o$ its
expiry, $n_o$ its signed contract count and $\mu_o$ its contract multiplier.
Write $\beta_o=n_o\mu_o$ and $\phi_o=+1$ for a call, $-1$ for a put.
Observed marks $V_{o0}^{\rm market}$ are per quoted unit. The multiplier is
applied exactly once, and all monetary inputs must use one currency.

With valuation date $t_0$ and horizon $t_h$, ACT/365 maturity is
\[
 \tau_o=\frac{\operatorname{days}(T_o-t_0)}{365},\qquad
 \tau'_o=\frac{\operatorname{days}(T_o-t_h)}{365}.
\]
The option must be unexpired at $t_0$. At $\tau'_o=0$ use
$[\phi_o(S-K_o)]_+$. Expiry before $t_h$ fails because the terminal scenario
does not supply settlement data. Underlying return observations must match the
horizon; changing the option date gap does not rescale the fitted return model.
Nominal mark-change P\&L excludes funding, stock-dividend cash credits and FX.

For aligned horizon-compatible historical log returns $\ell_{ti}$, oldest first,
use the same normalized exponential weights in all fitted moments:
\[
 \omega_t=\frac{\alpha^{n-t}}{\sum_u\alpha^{n-u}},\quad
 c_i=\sum_t\omega_t\ell_{ti},\quad
 \sigma_i^2=\sum_t\omega_t(\ell_{ti}-c_i)^2,\quad
 x_{ti}=\frac{\ell_{ti}-c_i}{\sigma_i}.
\]
For nonzero variance, let $D=\diag(\sigma_i)$ and
$C=\sum_t\omega_t x_tx_t^\top$, with descending eigenpairs $(u_a,\lambda_a)$.
The existing PCA layer owns degenerate-history validation and numerical scaling;
this extension does not change that layer. Retaining $p$ directions gives
\[
 \epsilon_t=D(I-U_pU_p^\top)x_t,\qquad
 R=\sum_t\omega_t\epsilon_t\epsilon_t^\top.
\]
At a smooth portfolio center, let $v$ be the full stock/option log-return
gradient. The baseline residual direction and factor matrix are
\[
 a_{\rm res}=\frac{Rv}{\sqrt{v^\top Rv}},\quad
 A=[Du_1\sqrt{\lambda_1},\ldots,Du_p\sqrt{\lambda_p},a_{\rm res}],
\]
with $a_{\rm res}=0$ if $v^\top Rv=0$. Weighted observation products avoid a dense
residual covariance. This preserves local linear residual variance; it can miss
residual gamma risk, particularly in delta-neutral option books.

All positions on underlying $i$ share
\[
 \ell(z)=c+Az,\qquad S_i(z)=S_{i0}e^{c_i+a_i^\top z}
 =\bar S_i e^{a_i^\top z},\qquad \|z\|_2\leq\rho.
\]
The configured radius is a stress budget. PCA alone does not give it a
probability interpretation. At $z=0$, log return is $c$, not zero.

\newpage
\textbf{2. European and American valuation}

For positive remaining maturity $\tau$, volatility $\nu$, rate $r$ and continuous
yield $q$, define
\[
 d_1=\frac{\log(S/K)+(r-q+\nu^2/2)\tau}{\nu\sqrt\tau},\quad
 d_2=d_1-\nu\sqrt\tau,
\]
\[
 V_E=\phi[Se^{-q\tau}\Phi(\phi d_1)-Ke^{-r\tau}\Phi(\phi d_2)],\quad
 \Delta_E=\phi e^{-q\tau}\Phi(\phi d_1),\quad
 \Gamma_E=\frac{e^{-q\tau}\varphi(d_1)}{S\nu\sqrt\tau}.
\]
European spot gamma is positive for both calls and puts. The maturity derivative
at fixed spot and parameters is
\[
 T_E=\frac{Se^{-q\tau}\varphi(d_1)\nu}{2\sqrt\tau}
 -\phi qSe^{-q\tau}\Phi(\phi d_1)+\phi rKe^{-r\tau}\Phi(\phi d_2).
\]

\textbf{Ju--Zhong context.}
For an American option, first use the European-equivalent shortcuts: $q=0$ for
a call, or $r=0$ for a put, under the nonnegative rate/yield restrictions.
Otherwise define
\[
 \beta_J=\frac{2(r-q)}{\nu^2},\qquad
 \kappa_J=\begin{cases}
 2r/[\nu^2(1-e^{-r\tau})],&r>0,\\
 2/(\nu^2\tau),&r=0,
 \end{cases}
\]
\[
 \zeta_J=\sqrt{(\beta_J-1)^2+4\kappa_J},\quad
 \eta_J=\frac{-(\beta_J-1)+\phi\zeta_J}{2},\quad
 d_J=\phi\zeta_J.
\]
The implementation uses \texttt{expm1} for $1-e^{-r\tau}$ and a stable quadratic
root computation. The finite boundary $S^\star$ solves
\[
 U^\star=\phi(S^\star-K)-V_E(S^\star),\qquad
 \Delta_E(S^\star)+\eta_J U^\star/S^\star-\phi=0.
\]
A sign-changing log-spot bracket and safeguarded scalar root solver determine
the boundary. Its residual and positive premium are checked before computing
\[
 b_J=-\frac{e^{-r\tau}\kappa_J^2}{2\zeta_J^2},\qquad
 c_J=-\frac{1}{d_J}\left[
 \frac{2T_E(S^\star)}{\nu^2U^\star}+e^{-r\tau}\kappa_J+2b_J\right].
\]
For $X=\log(S/S^\star)$, set $\mathcal D=1-b_JX^2-c_JX$ and
$\mathcal E=U^\star e^{\eta_JX}/\mathcal D$. The complete price is
\[
 \widehat V_J(S)=\begin{cases}
 V_E(S)+\mathcal E(X),&\phi(S^\star-S)>0,\\
 \phi(S-K),&\phi(S^\star-S)\leq0.
 \end{cases}
\]
The context is cached per contract and remaining maturity. Spot-only scenarios
do not recalibrate IV or solve another boundary. The final rational/exponential
pricing formula is not itself a quadratic polynomial.

\newpage
\textbf{3. Calibration, derivatives and pricing-domain checks}

\textbf{Mark inversion.}
The new book solves $V(\nu_o)=V_{o0}^{\rm market}$ using the contract's exercise
style. Brent's method maintains a bracket; every trial American IV rebuilds the
boundary. The default explicit bracket is $[0.01,5]$. Invalid trial contexts and
missing brackets fail. The final quoted-price residual must satisfy
\[
 |V(\nu_o)-V_{o0}^{\rm market}|\leq10^{-10}\max(1,V_{o0}^{\rm market}).
\]
Necessary economic mark bounds are checked, including the zero-volatility
American stopping lower limit
$\max_{0\leq u\leq\tau}[\phi(Se^{-qu}-Ke^{-ru})]_+$.
An American intrinsic-value plateau does not identify a unique positive IV.
Such a mark requires an explicitly supplied IV; supplied-IV status and its
calibration residual are recorded. Observed marks remain the P\&L anchor.
An American bracketed root is not a uniqueness certificate.

\textbf{Derivatives of the selected formula.}
Write $x=\log S$, $u_J=c_J+2b_JX$ and
$L_J=\eta_J+u_J/\mathcal D$. In continuation,
\begin{align*}
 D^J:=\widehat V_{J,x}
 &=S\Delta_E+\mathcal E L_J,\\
 G^J:=\widehat V_{J,xx}
 &=S\Delta_E+S^2\Gamma_E+
 \mathcal E\left[L_J^2+\frac{2b_J}{\mathcal D}
                      +\frac{u_J^2}{\mathcal D^2}\right].
\end{align*}
The European terms use the current spot. In the strict exercise region,
$D^J=G^J=\phi S$. At expiry, these derivatives apply inside the money, and both
are zero outside the money. Two-sided derivatives at the strike are undefined.

The final Ju--Zhong formula matches value at $S^\star$ but generally has a slope
jump:
\[
 \Delta_J^{\rm continuation}(S^\star)
 -\Delta_J^{\rm exercise}(S^\star)=\frac{c_JU^\star}{S^\star}.
\]
The implementation raises for requested derivatives within log-distance
$10^{-10}$ of the boundary; it does not invent a two-sided Hessian.
Price-only candidate evaluation remains valid.

\textbf{Domain checks.}
Each contract's marginal stress interval is
\[
 [S_i^{\min},S_i^{\max}]=
 [\bar S_i e^{-\rho\|a_i\|},\bar S_i e^{\rho\|a_i\|}].
\]
On its continuation portion, the exact quadratic denominator minimum is tested
at endpoints and the vertex. A value $\mathcal D\leq10^{-10}$ fails. Economic
bounds are sampled at 65 log-spaced points and checked at every evaluated
scenario: $\max(I,V_E)\leq\widehat V_J$, call value $\leq S$, put value $\leq K$,
with materiality tolerance $10^{-7}\max(S,K)$. These sampled checks are not an
interval-wide certificate. Nonfinite prices, invalid premiums and material
bounds violations fail explicitly. No clipping or scenario-dependent backend
switch is used; no automatic tree fallback is claimed.

\newpage
\textbf{4. Mixed P\&L, residual alignment and anchor polynomials}

Use $V_o=V_E$ for European contracts and $V_o=\widehat V_J$ for American
contracts, with the specified payoff/equivalent branches. Then
\[
 P(z)=\sum_i w_i\operatorname{expm1}(c_i+a_i^\top z)
 +\sum_o\beta_o[V_o(S_{i(o)}(z),\tau'_o)-V_{o0}^{\rm market}].
\]
This includes horizon time decay in both full P\&L and the Taylor constant.
For a smooth center, aggregate by underlying
\[
 v_i=w_i e^{c_i}+\sum_{o:i(o)=i}\beta_o V_{o,x}(\bar S_i),\qquad
 h_i=w_i e^{c_i}+\sum_{o:i(o)=i}\beta_o V_{o,xx}(\bar S_i).
\]
The full $v$, including American and European log deltas, selects the residual
direction. The center polynomial is
\[
 F(z)=p_0+g^\top z+\tfrac12z^\top Hz,\quad
 p_0=P(0),\quad g=A^\top v,\quad H=A^\top\diag(h)A.
\]
The Hessian is symmetrized for floating-point roundoff. Signed positions and
all legs on one underlying are aggregated before optimization; there is no
independent worst scenario for each option.

At a nonsmooth center, default PCA construction fails instead of manufacturing
Greeks. An explicit \texttt{localPnlGradient} may select alternative residual
alignment, or the caller may supply fixed geometry directly. The diagnostic
distinguishes mixed-delta, supplied and fixed-geometry construction. This choice
does not make the nonsmooth center differentiable. Price-only seeds and other
smooth anchors remain available. The optional extra negative-curvature residual
mode proposed in the supplied Markdown is not enabled by this implementation.

\textbf{Reachable boundaries.}
For an American boundary $B_o=S_o^\star$, or expiry strike $B_o=K_o$, put
$b_o^\star=\log(B_o/\bar S_{i(o)})$. The boundary plane is
\[
 a_{i(o)}^\top z=b_o^\star.
\]
For a nonzero loading row, it intersects the ball exactly when
$|b_o^\star|\leq\rho\|a_{i(o)}\|$, and its nearest point is
$z_o^\star=b_o^\star a_{i(o)}/\|a_{i(o)}\|^2$.
Zero loading is handled separately without division. A boundary touching the
ball or a coarse lattice need not have feasible points on both sides.

The search generates center and axis seeds, nearby boundary-side seeds and
marginal extreme seeds. After integer rounding/projection it checks which
boundary sides are actually represented. Extreme seeds help expose the loss
of a short expiry option whose center lies on a flat out-of-the-money payoff.

At each smooth lattice anchor $z_j$, form
\[
 F_j(z)=P(z_j)+g_j^\top(z-z_j)+\tfrac12(z-z_j)^\top H_j(z-z_j).
\]
In global coordinates its coefficients are
$\widetilde g_j=g_j-H_jz_j$ and
$\widetilde p_j=P(z_j)-g_j^\top z_j+\tfrac12z_j^\top H_jz_j$.
Keep $A,c,\rho$ fixed. Nonsmooth anchors are fully repriced but do not generate
derivative QUBOs. Multi-anchor global-ball proposals are heuristic; the optional
local trust-ball encoding from the Markdown is not included.

\newpage
\textbf{5. Unchanged integer encoding and nonlinear selection}

For $k\in\{2,\ldots,8\}$ bits per coordinate, write
\[
 m=2^{k-1}-1,\quad t_a=\sum_{j=0}^{k-1}2^j b_{aj}-m,\quad
 z_a=\frac{\rho}{m}t_a,\quad \sum_a t_a^2\leq m^2.
\]
The extra representable unsigned level is excluded by the exact ball budget.
For product bits $y_{aij}=b_{ai}b_{aj}$, the squared integer norm is represented
by the linear expression
\[
 Q(b,y)=dm^2+\sum_{a,j}(2^{2j}-2m2^j)b_{aj}
             +2\sum_a\sum_{i<j}2^{i+j}y_{aij}.
\]
Let $s=\sum_{h=0}^{L-1}2^h u_h$ with $L=\lceil\log_2(m^2+1)\rceil$.
For the actual anchor polynomial $F_j$, the QUBO is
\begin{align*}
 E(b,y,u)={}&F_j(z(b))+\mathcal P[Q(b,y)+s-m^2]^2\\
 &+\mathcal P\sum_a\sum_{i<j}
 [b_{ai}b_{aj}-2b_{ai}y_{aij}-2b_{aj}y_{aij}+3y_{aij}].
\end{align*}
Each product penalty is nonnegative and vanishes exactly for a consistent
product. An invalid integer encoding incurs an unweighted penalty of at least
one. If $B$ is the sum of absolute nonconstant binary objective coefficients,
the default $\mathcal P=1.1\max(B,\varepsilon)$ is sufficient for global
feasibility in exact arithmetic. Each anchor recomputes $B$ for its polynomial.
Weaker configured penalty multipliers may remove that guarantee.

The number of variables remains
\[
 N=dk+d\binom{k}{2}+L.
\]
For $d=3,k=8$, this is $24+84+14=122$. Option count increases pricing and
anchor-search cost, but not the variable count of each QUBO.

\textbf{Repair and ranking.}
Every raw solver output receives integer/product/slack diagnostics. The new
search projects its integer coordinates into the ball and rebuilds auxiliaries
before pricing; an invalid raw sample can lie outside the screened domain.
Deterministic seed candidates are also repaired. Nonlinear local improvement
prices only feasible neighbors and accepts strict decreases in complete $P$.
The original 26-neighbor mode applies at three dimensions, with explicit
pairwise repair above three dimensions. Convergence versus the move cap is
retained for every candidate.

For all retained candidates $\mathcal C$, the reported result is
\[
 M_{\rm found}=\max_{z\in\mathcal C}\max(0,-P(z)).
\]
The result includes the worst coordinates and stressed spots; stock, European
and American P\&L; repaired feasible encodings; raw feasibility; anchor counts;
and lattice boundary-side coverage. Calibration and pricing-domain diagnostics
are retained by the book/model and archived by the runnable example.
The existing solver interface returns one selected binary sample per solve;
this search ranks those returned samples and its deterministic seeds after
repair, not every internal solver trajectory.

No stock-only convexity certificate applies to a nonzero option position.
Neither QUBO optimality nor neighborhood convergence establishes a global
minimum of the nonlinear American portfolio value.

\newpage
\textbf{6. Validation, reproducibility and interpretation}

The implemented formula reproduces these supplied numerical examples:
\begin{center}\small
\begin{tabular}{p{.62\linewidth}r}
\toprule
Inputs ($S=K=100$) & Ju--Zhong price\\\midrule
Call: $\tau=0.5$, $r=0.03$, $q=0.07$, $\nu=0.2$ & 4.76818024\\
Put: $\tau=3$, $r=0.08$, $q=0$, $\nu=0.2$ & 6.95561233\\\bottomrule
\end{tabular}\end{center}
For the put, the implementation obtains
\[
 S^\star=82.2759889300,\quad U^\star=8.3733727010,\quad
 \Delta_J^{\rm cont}-\Delta_J^{\rm exercise}=0.03133987185.
\]
Independent log-price finite differences agree with analytic first/second
derivatives on smooth continuation and exercise branches. Existing independent
CRR trees at 800 and 1,600 steps refine to within 0.01 quoted units for these
two cases; the Ju--Zhong/reference difference remains visible. This validates
implementation examples, not a uniform American-pricing error bound.

The focused suites passed 48 tests. The complete Python suite ran 363 tests
successfully with 9 skips. Checks include mixed-style calibration, explicit-IV
residuals, zero-rate limits, immutable contexts, signed multipliers, calendar
maturity, mixed gradients/Hessians, residual alignment, exact QUBO energy,
fixed variable counts, nonsmooth anchors, expiry plateau escape, boundary-side
coverage and feasible-only repair. Native and GPU solver kernels are unchanged;
no separate GPU benchmark or native rebuild was needed for this CPU extension.

Run the synthetic end-to-end example from the repository root:
\begin{verbatim}
PYTHONPATH=src python tools/example_option_factor_stress.py
\end{verbatim}
To archive it, add \texttt{--output} followed by
\path{docs/benchmarks/american_factor_stress_example.json}.
It fits underlying PCA, creates European/American call and put marks, calibrates
IV, solves the anchor QUBOs with simulated annealing, repairs and selects the
greatest model loss. Its JSON separates same-IV CRR prices at the final worst
scenario from the model margin. It is an offline composition demonstration,
not a market-data coverage study or an exercise/assignment backtest.

\textbf{Separate approximation limits.}
Residual reduction, finite lattice resolution, Taylor proposals, finite solver
search and Ju--Zhong pricing each introduce different limitations. Nonlinear
repricing removes Taylor error from the reported value of a retained candidate;
it does not recover unsearched scenarios or remove the American-pricing
approximation. The additional curvature residual and local trust-ball designs
in the supplied Markdown remain optional future extensions. Domain failures
are explicit; a material unpriced contract is never assigned zero risk.

\textbf{Sources and implementation notes.}
The primary specification is the supplied \path{american_pca_qubo_model.md},
dated 12 September 2026. Its reported 87-case paper-table audit was not rerun
here; only the implementation tests described above are claimed. The cited
Ju--Zhong paper is N. Ju and R. Zhong, ``An Approximate Formula for Pricing
American Options,'' \emph{The Journal of Derivatives} 7(2), 1999, 31--40.
The stable coefficients and corrected European Greek terms were cross-checked
against the \href{https://github.com/lballabio/QuantLib/blob/master/ql/pricingengines/vanilla/juquadraticengine.cpp}{official QuantLib Ju engine source};
QuantLib is not a dependency. Further details and commands are in
\path{docs/benchmarks/american_factor_stress.md}.
\end{document}
